\VerdierPhysicalPageAnchor{0028}%
\VerdierSemanticAnchor{intro-7-locally-compact-model}
particularly simple case of locally compact spaces provides a model for the
statements one obtains. Let $X$ and $Y$ be two locally compact spaces (say,
of finite cohomological dimension), $A$ a (noetherian) ring,
$\mathbf D(X_A)$ and $\mathbf D(Y_A)$ the derived categories of the
categories $X_A$ and $Y_A$ of sheaves of $A$-modules on $X$ and $Y$
respectively, and $f\colon X\longrightarrow Y$ a continuous map. In this
situation one can then define:

\VerdierSemanticAnchor{intro-7-f-shriek}
\begin{enumerate}
\renewcommand{\labelenumi}{\alph{enumi})}
\item A functor
\[
f_!\colon X_A\longrightarrow Y_A,
\]
the direct image with proper support, and the corresponding total derived
functor:
\VerdierSemanticAnchor{intro-7-rf-shriek}
\[
\mathrm Rf_!\colon \mathbf D(X_A)\longrightarrow \mathbf D(Y_A).
\]
(When $f\colon X\longrightarrow Y$ is a proper map, $f_!$ is the usual
direct-image functor.)

\VerdierSemanticAnchor{intro-7-f-upper-shriek}
\item A functor
\[
f^!\colon \mathbf D^+(Y_A)\longrightarrow \mathbf D^+(X_A)
\]
(the sign $+$ means that one restricts to complexes bounded below), whose
precise definition is fairly technical but which is characterized, up to
unique isomorphism, by property c).

\VerdierSemanticAnchor{intro-7-duality-adjunction}
\item A bifunctorial isomorphism:
\VerdierSemanticAnchor{intro-7-delta-f}
\[
\Delta_f(F,G)\colon
\operatorname{RHom}(\mathrm Rf_!F,G)
\longrightarrow
\operatorname{RHom}(F,f^!G)
\]
defined for every object $F$ of $\mathbf D(X_A)$ and every object $G$ of
$\mathbf D^+(Y_A)$ ($\operatorname{RHom}$ denotes the total derived functor
of the “complex of homomorphisms” functor).
\end{enumerate}

The important phenomena are then the following:

\VerdierSemanticAnchor{intro-7-duality-D1}
\noindent\textit{D1)} Formation of the functor $f^!$ is transitive as $f$
varies.

\VerdierSemanticAnchor{intro-7-duality-D2}
\noindent\textit{D2)} The isomorphism $\Delta_f$ is compatible with
composition of maps.

\VerdierSemanticAnchor{intro-7-duality-D3}
\noindent\textit{D3)} When $X$ is smooth over $Y$ of relative dimension
$n$, i.e. when, locally on $X$, $X$ is $Y$-isomorphic to the product of an
open subset of $Y$ with the topological space $\mathbf R^n$ (for example:
$Y$ is a point and $X$ a topological manifold), then for every object $G$ of
$\mathbf D^+(Y_A)$ there is a canonical isomorphism
\VerdierSemanticAnchor{intro-7-relative-orientation}
\[
f^!(G)\simeq f^*(G)\otimes\omega_{X/Y}[n],
\]
where $f^*(G)$ is the usual inverse image, $\omega_{X/Y}$ is a locally free
sheaf of rank $1$, namely the relative orientation sheaf of $X$ over $Y$,
and $\omega_{X/Y}[n]$ denotes the complex obtained by shifting the degrees of
$\omega_{X/Y}$ by $n$ units.
